\documentclass[12pt, a4paper]{article}

% ── Packages ──────────────────────────────────────────────────────────────────
\usepackage[a4paper, margin=2.5cm]{geometry}
\usepackage{graphicx}
\usepackage{float}
\usepackage{amsmath}
\usepackage{amssymb}
\usepackage{xcolor}
\usepackage{mdframed}
\usepackage{parskip}
\usepackage{enumitem}
\usepackage{titlesec}
\usepackage{booktabs}
\usepackage{hyperref}
\usepackage[T1]{fontenc}
\usepackage{tikz}
\usepackage{listings}
\usetikzlibrary{arrows.meta}

% ── Graphics path ─────────────────────────────────────────────────────────────
\graphicspath{{../../Images_for_latex/Lecture_4/}}

% ── Colours ───────────────────────────────────────────────────────────────────
\definecolor{keyblue}{RGB}{30, 80, 160}
\definecolor{keybluebg}{RGB}{220, 235, 255}
\definecolor{noteorange}{RGB}{200, 100, 0}
\definecolor{noteorangebg}{RGB}{255, 240, 215}
\definecolor{warnred}{RGB}{180, 30, 30}
\definecolor{warnredbg}{RGB}{255, 220, 220}
\definecolor{defgreen}{RGB}{20, 110, 50}
\definecolor{defgreenbg}{RGB}{215, 245, 225}
\definecolor{headernavy}{RGB}{25, 35, 90}

% ── Box environments ──────────────────────────────────────────────────────────
\newmdenv[linecolor=keyblue,backgroundcolor=keybluebg,linewidth=1.5pt,
  leftline=true,rightline=false,topline=false,bottomline=false,
  innerleftmargin=10pt,innerrightmargin=8pt,
  innertopmargin=6pt,innerbottommargin=6pt,
  skipabove=8pt,skipbelow=8pt]{keybox}

\newmdenv[linecolor=noteorange,backgroundcolor=noteorangebg,linewidth=1.5pt,
  leftline=true,rightline=false,topline=false,bottomline=false,
  innerleftmargin=10pt,innerrightmargin=8pt,
  innertopmargin=6pt,innerbottommargin=6pt,
  skipabove=8pt,skipbelow=8pt]{notebox}

\newmdenv[linecolor=warnred,backgroundcolor=warnredbg,linewidth=1.5pt,
  leftline=true,rightline=false,topline=false,bottomline=false,
  innerleftmargin=10pt,innerrightmargin=8pt,
  innertopmargin=6pt,innerbottommargin=6pt,
  skipabove=8pt,skipbelow=8pt]{warnbox}

\newmdenv[linecolor=defgreen,backgroundcolor=defgreenbg,linewidth=1.5pt,
  leftline=true,rightline=false,topline=false,bottomline=false,
  innerleftmargin=10pt,innerrightmargin=8pt,
  innertopmargin=6pt,innerbottommargin=6pt,
  skipabove=8pt,skipbelow=8pt]{defbox}

% ── Section formatting ────────────────────────────────────────────────────────
\titleformat{\section}[block]{\large\bfseries\color{headernavy}}{}{0pt}{}
  [\vspace{2pt}\rule{\linewidth}{0.8pt}\vspace{4pt}]
\titleformat{\subsection}[block]{\normalsize\bfseries\color{keyblue}}{}{0pt}{}

% ── Document ──────────────────────────────────────────────────────────────────
\begin{document}
\section*{Exercises: We do All 4 Both only present the 4-6}
\vspace{-0.5cm}
\begin{enumerate}
  \setcounter{enumi}{2}
  \item Write a Python function that determines if a straight line intersects
        a specified circle, returning the intersection coordinates.
  \item Use the function above to emulate a LiDAR shooting at a cylindrical
        robot as in Wasik et al. Plot the robot and the point cloud.
        Add noise to the measurements.
  \item Place the simulated robot in front of a wall.
        Simulate the point cloud from both robot and wall.
  \item Assign wall point cloud data to one cluster. Use the optimisation
        from Ex.\ 2 to determine the equation of the corresponding line.
\end{enumerate}
\vspace{-0.2cm}
\subsection*{Exercise 4: LiDAR point cloud on cylindrical robot}
\vspace{-0.2cm}
For each beam angle $\alpha_i \in [0, 2\pi]$ we check whether the ray from the
LiDAR origin hits the cylindrical robot modelled as a circle at position
$(x_r, y_r)$ with radius $r$. The ray direction is:
\[
  (dx, dy) = (\cos\alpha_i,\; \sin\alpha_i)
\]
The intersection is checked by substituting the parametric ray
$P = \text{lidar\_pos} + t\cdot(dx,\,dy)$ into the circle boundary ---
a point is on the circle when its distance to the center equals the radius,
so we require $\|P - \text{robot\_pos}\| = r$, or equivalently:
\[
  \|P - \text{robot\_pos}\|^2 = r^2
\]
This gives one equation in one unknown $t$ --- solve for $t$, then plug back
in to get the hit point. Returns the closest hit $(x_h, y_h)$
if $t > 0$ and within \texttt{lidar\_range}, otherwise \texttt{None}.

\begin{notebox}
Substituting the ray $P = \text{lidar\_pos} + t\cdot(dx,\,dy)$ into the
circle boundary condition:
\[
  \|P - \text{robot\_pos}\|^2 = r^2 \quad \rightarrow \quad \|P - \text{robot\_pos}\|^2 + r^2 = 0
\]
expands to a quadratic in $t$:
\[
  \underbrace{(dx^2+dy^2)}_{a_c}\,t^2
  + \underbrace{2(fx\cdot dx + fy\cdot dy)}_{b_c}\,t
  + \underbrace{(fx^2+fy^2-r^2)}_{c_c} = 0
\]
where $(fx,\,fy) = \text{lidar\_pos} - \text{robot\_pos}$.
Solving with the quadratic formula:
\[
  t = \frac{-b_c - \sqrt{b_c^2 - 4a_c c_c}}{2a_c}
\]
The smallest positive $t$ gives the closest hit point:
\[
  (x_h,\,y_h) = \text{lidar\_pos} + t\cdot(dx,\,dy)
\]
\end{notebox}
\textbf{Noise model.} Each hit point is converted to polar coordinates
relative to the LiDAR origin --- taking the difference in $x$ and $y$
between hit point and lidar gives the vector between them,
Pythagoras gives the true range and \texttt{atan2} gives the true bearing:
\[
  d = \sqrt{(x_r - x_l)^2 + (y_r - y_l)^2}, \qquad
  \alpha = \text{atan2}(y_r - y_l,\; x_r - x_l)
\]
where $x_l, y_l$ is the lidar position and $x_r, y_r$ the hit point on the robot.

Independent Gaussian noise is then added in polar domain:
\[
  \tilde{d} = d + \nu_d, \quad \nu_d \sim \mathcal{N}(0, \sigma_d^2)
  \qquad
  \tilde{\alpha} = \alpha + \nu_\alpha, \quad \nu_\alpha \sim \mathcal{N}(0, \sigma_\alpha^2)
\]
and converted back to Cartesian:
\[
  \tilde{x}_r = x_l + \tilde{d}\cos(\tilde{\alpha}), \qquad
  \tilde{y}_r = y_l + \tilde{d}\sin(\tilde{\alpha})
\]
Noise is added in polar domain because angle noise scales with distance ---
a $1^\circ$ error at distance $d$ causes a lateral shift of
$d\cdot\sin(\nu_\alpha) \approx d\cdot\nu_\alpha$,
while range noise $\sigma_d$ is constant regardless of distance.

\begin{warnbox}
This is precisely the \textbf{uncertainty ellipse} --- we are more certain
about the range (radial axis, fixed half-length $\approx \sigma_d$) than
about the lateral position (tangential axis, half-length $\approx d\cdot\sigma_\alpha$),
because the same angle error $\sigma_\alpha$ produces a larger sideways
displacement the further the object is --- $\Delta = d\cdot\sigma_\alpha$
grows with distance while the range error stays constant.
\end{warnbox}
\vspace{-0.5cm}
\begin{notebox}
  \textbf{Why angle noise is more dangerous at long range.}
  A $1^\circ$ angle error causes a lateral shift of $\Delta = d\cdot\sin(\nu_\alpha)$:
  \[
    d = 1\,\text{m}:  \quad \Delta = 1 \cdot \sin(1^\circ) \approx 1 \cdot 0.017 = 1.7\,\text{cm}
  \]
  \[
    d = 10\,\text{m}: \quad \Delta = 10 \cdot \sin(1^\circ) \approx 10 \cdot 0.017 = 17\,\text{cm}
  \]
  Range noise $\sigma_d$ is constant: $5\,\text{cm}$ error at $1$\,m is the
  same $5\,\text{cm}$ at $10$\,m. So angular error scales linearly with distance --- the further the object,
  the larger the lateral shift from the same angle mistake.
\end{notebox}


\vspace{-0.6cm}
\subsection*{Exercise 5: Robot in front of wall}
\vspace{-0.3cm}
We add a wall behind the robot. For each beam angle we check the robot first
--- if it hits, the wall behind is occluded:
\[
  \text{for each } \alpha_i: \quad
  \begin{cases}
    \text{check robot} \rightarrow \text{hit} & \rightarrow \text{record robot point, skip wall}\\
    \text{check robot} \rightarrow \text{miss} & \rightarrow \text{check wall}
  \end{cases}
\]
The wall intersection uses two parametric lines --- the ray and the wall segment:
\[
  \text{ray:} \quad P = \text{lidar} + t\cdot(dx,\,dy)
  \qquad
  \text{wall:} \quad W = \text{wall\_start} + s\cdot(wx,\,wy)
\]
Setting $P = W$ and solving gives $t$ and $s$ via the 2D cross product:
\[
  \text{denom} = dx\cdot wy - dy\cdot wx
\]
If $\text{denom} = 0$ the ray is parallel to the wall (cross product of two
parallel vectors is zero) --- no intersection possible. 

Otherwise we solve:
\[
  t = \frac{fx\cdot wy - fy\cdot wx}{\text{denom}}, \qquad
  s = \frac{fx\cdot dy - fy\cdot dx}{\text{denom}}
\]
where $(fx,\,fy) = \text{wall\_start} - \text{lidar}$.
Two conditions must both hold for a valid hit:
\begin{itemize}[leftmargin=1.4em]
  \item $0 < t \leq \texttt{lidar\_range}$ --- intersection is in front of lidar and within range
  \item $s \in [0,\,1]$ --- intersection lies within the actual wall segment, not on the infinite line
\end{itemize}

\begin{notebox}
  \textbf{Concrete example:} lidar at $(0,0)$, ray going right $(dx,dy)=(1,0)$,
  wall from $(5,-3)$ to $(5,3)$ so $(wx,wy)=(0,6)$.

  \textbf{Finding $t$ --- how far along the ray:}
  \[
  \begin{aligned}
    \text{denom} &= 1\cdot6 - 0\cdot0 = 6 \\
    fx &= 5,\quad fy = -3\\
    t &= \frac{5\cdot6 - (-3)\cdot0}{6} = 5
    \quad\rightarrow\quad \text{hit at }(0,0)+5\cdot(1,0) = (5,0) \;\checkmark
  \end{aligned}
  \]
  \textbf{Finding $s$ --- where on the wall segment:}
  \[
    s = \frac{5\cdot0 - (-3)\cdot1}{6} = 0.5
    \quad\rightarrow\quad \text{halfway along wall} \;\checkmark
  \]
  If hit at $(5,-4)$: $s < 0$ $\rightarrow$ miss.\quad
  If hit at $(5,\phantom{-}4)$: $s > 1$ $\rightarrow$ miss.
\end{notebox}
\subsection*{Exercise 6: Wall cluster and line fitting}

We want to fit a line to the wall point cloud by minimising the sum of
squared perpendicular distances from each point to the line $ax+by+c=0$:
\[
  \min_{a,b,c} \sum_{i} \frac{(ax_i + by_i + c)^2}{a^2 + b^2}
\]
where $(a,b)$ is the normal vector defining the line orientation and $c$
the offset from the origin. The denominator $a^2+b^2$ normalises for the
length of the normal vector so the result is a true perpendicular distance.
The optimisation is solved numerically using \texttt{scipy.optimize.minimize}.

The normal vector $(a, b)$ defines the orientation of the line,
and $c$ the perpendicular distance from the origin:
\begin{itemize}[leftmargin=1.4em]
  \item $a=1,\; b=0$ $\;\rightarrow\;$ normal points in $x$-direction
        $\;\rightarrow\;$ line is \textbf{vertical}
  \item $a=0,\; b=1$ $\;\rightarrow\;$ normal points in $y$-direction
        $\;\rightarrow\;$ line is \textbf{horizontal}
\end{itemize}

Once $a,b,c$ are found the line is reconstructed for plotting by solving
for $x$ from the line equation:
\[
  ax + by + c = 0 \quad\rightarrow\quad x = \frac{-(by + c)}{a}
\]
Describes how $x$ changes as $y$ moves along the line ---
for a perfect vertical wall $b \approx 0$ so $x \approx -c/a$ is constant,
confirming the line is vertical.
Evaluated at $y = [y_{\min},\; y_{\max}]$ of the wall points ---
constraining the plotted line to the actual wall segment.

\begin{notebox}
  \textbf{Example --- vertical wall at $x=5$:}
  $a=1,\; b=0$ because the wall is perpendicular to the $x$-axis,
  $c=-5$ because the origin is 5 units from the wall:
  \[
    \underbrace{a\cdot x}_{\text{how far in }x} +
    \underbrace{b\cdot y}_{\text{how far in }y} +
    \underbrace{c}_{\text{offset from origin}} = 0
    \quad\rightarrow\quad
    1\cdot5 + 0\cdot y + (-5) = 0 \;\checkmark
  \]
\end{notebox}

\end{document}